El copernici $^{277}_{112}\mathrm{Cn}$ va ser sintetitzat al laboratori del Centre per a la Recerca d'Ions Pesants (GSI) de Darmstadt (Alemanya) el 9 de febrer del 1999. El nom oficial data del febrer del 2010, en honor de Nicolau Copèrnic. Per a obtenir-lo, es bombardeja una diana de plom amb projectils d'àtoms de zinc. La reacció es pot escriure així:
\[ {}^{208}_{\ a}\mathrm{Pb} + {}^{70}_{\,b}\mathrm{Zn} \rightarrow {}^{277}_{112}\mathrm{Cn} + {?} \]

\figura[0.25]{fig1.png}

El $^{277}_{112}\mathrm{Cn}$ es desintegra segons la seqüència següent:
\begin{align*}
&{}^{277}_{112}\mathrm{Cn} \rightarrow {}^{273}_{110}\mathrm{X} + {?}\\
&{}^{273}_{110}\mathrm{X} \rightarrow {}^{269}_{108}\mathrm{X} + {?}\\
&{}^{269}_{108}\mathrm{X} \rightarrow {}^{265}_{106}\mathrm{X} + {?}\\
&{}^{265}_{106}\mathrm{X} \rightarrow {}^{261}_{104}\mathrm{X} + {?}\\
&{}^{261}_{104}\mathrm{X} \rightarrow {}^{257}_{102}\mathrm{X} + {?}\\
&{}^{257}_{102}\mathrm{X} \rightarrow {}^{253}_{100}\mathrm{Fm} + {?}
\end{align*}

El $^{277}_{112}\mathrm{Cn}$ té un període de semidesintegració de 0,17~ms.

\begin{apartat}{1}
Completeu la reacció d'obtenció del $^{277}_{112}\mathrm{Cn}$ a partir de plom i de zinc. Quin tant per cent de $^{277}_{112}\mathrm{Cn}$ roman sense desintegrar-se al cap d'un minut d'haver-se produït la reacció d'obtenció d'aquest isòtop?
\end{apartat}

\begin{apartat}{1}
Escriviu la seqüència o sèrie radioactiva (amb tots els símbols dels elements) fins a arribar al fermi.
\end{apartat}

\dades{%

  \begin{tabular}{|c|c|c|c|c|c|c|c|}
  \hline
  $_{82}\mathrm{Pb}$ & $_{110}\mathrm{Ds}$ & $_{108}\mathrm{Hs}$ & $_{106}\mathrm{Sg}$ & $_{104}\mathrm{Rf}$ & $_{102}\mathrm{No}$ & $_{100}\mathrm{Fm}$ & $_{30}\mathrm{Zn}$\\
  \hline
  plom & darmstadti & hassi & seaborgi & rutherfordi & nobeli & fermi & zinc\\
  \hline
  \end{tabular}}
